% ============================================================
%  Shared preamble: USABO, IBO and Beyond -- Biology Compendium
% ============================================================
\documentclass[10pt,letterpaper]{article}

\usepackage[T1]{fontenc}
\usepackage[utf8]{inputenc}
\usepackage[letterpaper,left=1.10in,right=1.10in,top=1.00in,bottom=0.95in,
            headsep=17pt,footskip=28pt]{geometry}

% ---- Fonts -------------------------------------------------
\usepackage{mathpazo}
\usepackage[scaled=0.90]{helvet}
\usepackage{courier}
\usepackage[expansion=false]{microtype}
\linespread{1.22}
\setlength{\parskip}{8.5pt}
\setlength{\parindent}{0pt}
\setlength{\emergencystretch}{3em}

% ---- Mathematics and chemistry -----------------------------
\usepackage{amsmath,amssymb,amsthm,mathtools}
\usepackage{bm}
\usepackage[version=4]{mhchem}
\allowdisplaybreaks

% ---- Graphics and tables -----------------------------------
\usepackage{graphicx}
\usepackage{booktabs}
\usepackage{array}
\usepackage{multicol}
\usepackage{adjustbox}
\usepackage{longtable}
\usepackage{enumitem}
\setlist{nosep,topsep=3pt,itemsep=3pt,parsep=0pt,leftmargin=15pt}

\usepackage{tikz}
\usetikzlibrary{calc,intersections,angles,quotes,patterns,arrows.meta,positioning,decorations.pathreplacing}

% ---- Colours -----------------------------------------------
\usepackage{xcolor}
\definecolor{bioink}{HTML}{16332B}
\definecolor{biogreen}{HTML}{2C6E54}
\definecolor{bioline}{HTML}{C5D2CB}
\definecolor{biogrey}{HTML}{F3F6F4}
\definecolor{biogold}{HTML}{9C6B1E}
\definecolor{bioteal}{HTML}{1F6F6B}
\definecolor{biored}{HTML}{9B2C2C}
\definecolor{bioplum}{HTML}{6B3A63}
\definecolor{bionum}{HTML}{9BA8A1}
\definecolor{tiercore}{HTML}{2C6E54}
\definecolor{tierbeyond}{HTML}{6B5B3E}
\definecolor{tiercat}{HTML}{3A5A7A}

% ---- Boxes -------------------------------------------------
\usepackage[most]{tcolorbox}
\tcbset{biobase/.style={enhanced,arc=0pt,outer arc=0pt,
  left=11pt,right=0pt,top=5pt,bottom=5pt,
  boxrule=0pt,colframe=white,colback=white,
  fonttitle=\sffamily\bfseries\footnotesize,
  toptitle=0pt,bottomtitle=3pt,
  before skip=13pt,after skip=13pt}}

\newtcolorbox{keybox}[1][]{biobase,breakable,colback=biogrey,colbacktitle=biogrey,
  coltitle=biogreen,borderline west={2.2pt}{0pt}{biogreen},#1}
\newtcolorbox{thmbox}[1][]{biobase,breakable,colback=white,colbacktitle=white,
  coltitle=bioink,borderline west={2.2pt}{0pt}{bioink},#1}
\newtcolorbox{defbox}[1][]{biobase,breakable,colback=white,colbacktitle=white,
  coltitle=bioteal,borderline west={2.2pt}{0pt}{bioteal},#1}
\newtcolorbox{warnbox}[1][]{biobase,breakable,colback=white,colbacktitle=white,
  coltitle=biored,borderline west={2.2pt}{0pt}{biored},#1}
\newtcolorbox{tipbox}[1][]{biobase,breakable,colback=white,colbacktitle=white,
  coltitle=biogreen,borderline west={2.2pt}{0pt}{biogreen},#1}
\newtcolorbox{obsbox}[1][]{biobase,breakable,colback=white,colbacktitle=white,
  coltitle=biogold,borderline west={2.2pt}{0pt}{biogold},#1}
\newtcolorbox{calcbox}[1][]{biobase,breakable,colback=biogrey,colbacktitle=biogrey,
  coltitle=bioink,borderline west={2.2pt}{0pt}{bionum},#1}
\newtcolorbox{expbox}[1][]{biobase,breakable,colback=white,colbacktitle=white,
  coltitle=bioplum,borderline west={2.2pt}{0pt}{bioplum},#1}

% Non-breakable clones, for use inside minipages and tables
\newtcolorbox{keyboxnb}[1][]{biobase,colback=biogrey,colbacktitle=biogrey,
  coltitle=biogreen,borderline west={2.2pt}{0pt}{biogreen},#1}
\newtcolorbox{thmboxnb}[1][]{biobase,colback=white,colbacktitle=white,
  coltitle=bioink,borderline west={2.2pt}{0pt}{bioink},#1}
\newtcolorbox{calcboxnb}[1][]{biobase,colback=biogrey,colbacktitle=biogrey,
  coltitle=bioink,borderline west={2.2pt}{0pt}{bionum},#1}

% ---- Index -------------------------------------------------
\usepackage{imakeidx}

% ---- Sectioning --------------------------------------------
\usepackage{titlesec}
\usepackage{titletoc}
\usepackage[hidelinks]{hyperref}
\hypersetup{colorlinks=true,linkcolor=biogreen,urlcolor=biogreen,citecolor=biogreen,
  bookmarksnumbered=true}

\titleformat{\section}{\sffamily\bfseries\large\color{bioink}}{\thesection}{0.6em}{}
\titleformat{\subsection}{\sffamily\bfseries\normalsize\color{biogreen}}{\thesubsection}{0.6em}{}
\titlespacing*{\section}{0pt}{16pt}{6pt}
\titlespacing*{\subsection}{0pt}{12pt}{4pt}

% ---- Running heads -----------------------------------------
\usepackage{fancyhdr}
\pagestyle{fancy}
\fancyhf{}
\renewcommand{\headrulewidth}{0pt}
\fancyhead[L]{\sffamily\footnotesize\color{bionum}\leftmark}
\fancyfoot[C]{\sffamily\footnotesize\color{bionum}\thepage}

% ---- Small utilities ---------------------------------------
\newcommand{\tm}[1]{\textbf{#1}}
\newcommand{\dg}{^{\circ}}
\newcommand{\degc}{\,^{\circ}\mathrm{C}}
\newcommand{\sn}[1]{\textit{#1}}
\newcommand{\ttitle}[1]{\textbf{\textcolor{bioink}{#1}}}
\newcommand{\secchip}[1]{%
  \noindent{\sffamily\bfseries\fontsize{13.5}{16.5}\selectfont\color{bioink}\raggedright #1\par}%
  \vspace{4pt}%
  \noindent{\color{biogreen}\rule{\textwidth}{1.1pt}}\par}

% Named result: index entry plus a label the catalogue can point at
\newcommand{\nr}[2]{\index{#2}\label{nr:#1}}
\newcommand{\nrrow}[2]{#2\dotfill\pageref{nr:#1}\\}

% ---- Parts -------------------------------------------------
\newcounter{biopart}
\newcommand{\prepart}{}
\newcommand{\pref}[1]{Part~\ref{part:#1}}
\newcommand{\prefn}[1]{\ref{part:#1}}

\newcommand{\parttitle}[3][]{%
  \clearpage
  \refstepcounter{biopart}%
  \ifx\relax#1\relax\else\label{part:#1}\fi
  \markboth{#2}{#2}%
  \thispagestyle{fancy}%
  \vspace*{6pt}%
  \noindent{\color{bioline}\rule{\textwidth}{1pt}}\\[6pt]
  \noindent{\sffamily\bfseries\footnotesize\color{bionum}PART \thebiopart}\\[8pt]
  \noindent\begin{minipage}{\textwidth}\raggedright
    {\sffamily\bfseries\fontsize{22}{26}\selectfont\color{bioink}%
     \hyphenchar\font=-1\hyphenpenalty=10000\exhyphenpenalty=10000 #2\par}%
  \end{minipage}\\[10pt]
  \noindent{\color{biogreen}\rule{\textwidth}{1.4pt}}\\[8pt]
  \noindent\begin{minipage}{\textwidth}\raggedright
    {\fontsize{10.5}{14}\selectfont\color{biogreen} #3\par}%
  \end{minipage}
  \vspace{10pt}\par
  \addcontentsline{toc}{section}{\protect\numberline{\thebiopart}#2}%
}

% ---- Tier dividers -----------------------------------------
\newcommand{\tierpartline}[1]{\noindent\hangindent=12pt {\color{biogreen}\textbullet}\ #1\par\vspace{1pt}}
\newcommand{\tierdivider}[4]{%
  \clearpage
  \thispagestyle{empty}%
  \vspace*{0.08\textheight}
  \noindent{\color{tier#1}\rule{\textwidth}{5pt}}\\[3pt]
  \noindent{\color{bioline}\rule{\textwidth}{1pt}}\\[26pt]
  \noindent\begin{minipage}{\textwidth}\raggedright
    {\sffamily\bfseries\fontsize{30}{36}\selectfont\color{bioink} #2\par}%
  \end{minipage}\\[16pt]
  \noindent{\color{bioline}\rule{\textwidth}{1pt}}\\[12pt]
  \noindent\begin{minipage}{\textwidth}\raggedright
    {\fontsize{11.5}{16}\selectfont\color{biogreen} #3\par}%
  \end{minipage}\\[20pt]
  {\footnotesize\begin{multicols}{2}#4\end{multicols}}
  \vfill
  \clearpage
}

% ---- Table safety: never let a tabular run past the measure --
\usepackage{etoolbox}
\BeforeBeginEnvironment{tabular}{\begin{adjustbox}{max width=\linewidth}}
\AfterEndEnvironment{tabular}{\end{adjustbox}}

% Part display names, so a cross reference into a part that is not in this volume degrades
% to the part's name in italics instead of a dangling question mark.
\makeatletter
\expandafter\def\csname pname@foundations\endcsname{Foundations, Scale, and the Language of Biology}
\expandafter\def\csname pname@chemistry\endcsname{The Chemistry of Life}
\expandafter\def\csname pname@macromolecules\endcsname{Biological Macromolecules}
\expandafter\def\csname pname@enzymes\endcsname{Enzymes, Energetics, and Metabolic Control}
\expandafter\def\csname pname@cellstructure\endcsname{Cell Structure and the Organelles}
\expandafter\def\csname pname@membranes\endcsname{Membranes and Membrane Transport}
\expandafter\def\csname pname@respiration\endcsname{Cellular Respiration and Fermentation}
\expandafter\def\csname pname@photosynthesis\endcsname{Photosynthesis}
\expandafter\def\csname pname@signaling\endcsname{Cell Signaling and Signal Transduction}
\expandafter\def\csname pname@cytoskeleton\endcsname{The Cytoskeleton, Motility, and the Extracellular Matrix}
\expandafter\def\csname pname@cellcycle\endcsname{The Cell Cycle, Mitosis, and Meiosis}
\expandafter\def\csname pname@cancer\endcsname{Cancer and the Control of Cell Number}
\expandafter\def\csname pname@replication\endcsname{DNA Replication, Repair, and Chromosome Structure}
\expandafter\def\csname pname@transcription\endcsname{Transcription, RNA Processing, and the Genetic Code}
\expandafter\def\csname pname@translation\endcsname{Translation and Protein Targeting}
\expandafter\def\csname pname@generegulation\endcsname{Regulation of Gene Expression and Epigenetics}
\expandafter\def\csname pname@biotech\endcsname{Biotechnology and Molecular Laboratory Methods}
\expandafter\def\csname pname@mendelian\endcsname{Mendelian Genetics, Pedigrees, and Probability}
\expandafter\def\csname pname@linkage\endcsname{Linkage, Mapping, and Chromosomal Genetics}
\expandafter\def\csname pname@popgen\endcsname{Population Genetics and Hardy-Weinberg}
\expandafter\def\csname pname@evolution\endcsname{Evolutionary Mechanisms and Natural Selection}
\expandafter\def\csname pname@macroevolution\endcsname{Speciation, Macroevolution, and the History of Life}
\expandafter\def\csname pname@phylogenetics\endcsname{Phylogenetics and Biosystematics}
\expandafter\def\csname pname@prokaryotes\endcsname{Bacteria, Archaea, and the Prokaryotic World}
\expandafter\def\csname pname@viruses\endcsname{Viruses and Subcellular Agents}
\expandafter\def\csname pname@protists\endcsname{Protists and the Origin of Eukaryotes}
\expandafter\def\csname pname@fungi\endcsname{Fungi}
\expandafter\def\csname pname@plantdiversity\endcsname{Plant Diversity, Life Cycles, and Reproduction}
\expandafter\def\csname pname@plantanatomy\endcsname{Plant Anatomy, Tissues, and Growth}
\expandafter\def\csname pname@planttransport\endcsname{Plant Transport, Nutrition, and Water Relations}
\expandafter\def\csname pname@planthormones\endcsname{Plant Hormones, Signals, and Development}
\expandafter\def\csname pname@animaldiversity\endcsname{Animal Diversity and Body Plans}
\expandafter\def\csname pname@musculoskeletal\endcsname{Animal Tissues, Skeleton, Muscle, and Movement}
\expandafter\def\csname pname@digestion\endcsname{Animal Nutrition and Digestion}
\expandafter\def\csname pname@gasexchange\endcsname{Gas Exchange and Respiratory Systems}
\expandafter\def\csname pname@circulation\endcsname{Circulation, Blood, and the Heart}
\expandafter\def\csname pname@immunology\endcsname{Immunology}
\expandafter\def\csname pname@osmoregulation\endcsname{Osmoregulation, Excretion, and Thermoregulation}
\expandafter\def\csname pname@nervous\endcsname{The Nervous System and Neurophysiology}
\expandafter\def\csname pname@sensory\endcsname{Sensory Biology}
\expandafter\def\csname pname@endocrine\endcsname{The Endocrine System}
\expandafter\def\csname pname@development\endcsname{Animal Reproduction and Development}
\expandafter\def\csname pname@ethology\endcsname{Ethology and Animal Behavior}
\expandafter\def\csname pname@popecology\endcsname{Population and Community Ecology}
\expandafter\def\csname pname@ecosystems\endcsname{Ecosystems, Biomes, Biogeochemistry, and Conservation}
\expandafter\def\csname pname@examcraft\endcsname{Exam Craft: The Open Exam, the Semifinal, and the IBO}
\expandafter\def\csname pname@devbio\endcsname{Developmental Biology in Depth}
\expandafter\def\csname pname@neuro\endcsname{Neurobiology in Depth}
\expandafter\def\csname pname@immuno\endcsname{Immunology in Depth}
\expandafter\def\csname pname@microbio\endcsname{Microbiology and Virology in Depth}
\expandafter\def\csname pname@compphys\endcsname{Comparative and Clinical Physiology}
\expandafter\def\csname pname@structural\endcsname{Structural Biology and Protein Science}
\expandafter\def\csname pname@systems\endcsname{Systems Biology and Gene Networks}
\expandafter\def\csname pname@techniques\endcsname{Molecular Techniques in Depth}
\expandafter\def\csname pname@evogenomics\endcsname{Evolutionary Genomics and Molecular Evolution}
\expandafter\def\csname pname@practical\endcsname{The Practical Exam: Lab Skills, Histology, and Data Interpretation}
\expandafter\def\csname pname@f-units\endcsname{Units, Scaling, Dimensional Analysis, and Magnitudes}
\expandafter\def\csname pname@f-biostats\endcsname{Biostatistics: Distributions, Estimation, and Error}
\expandafter\def\csname pname@f-tests\endcsname{Hypothesis Tests, Chi-Square, and Goodness of Fit}
\expandafter\def\csname pname@f-probability\endcsname{Probability, Pedigrees, and Bayesian Genetics}
\expandafter\def\csname pname@f-mendelian\endcsname{Mendelian and Linkage Calculations}
\expandafter\def\csname pname@f-hardyweinberg\endcsname{Hardy-Weinberg and Population Genetics Equations}
\expandafter\def\csname pname@f-selection\endcsname{Selection, Drift, Mutation, Migration, and Inbreeding}
\expandafter\def\csname pname@f-kinetics\endcsname{Enzyme Kinetics Equations}
\expandafter\def\csname pname@f-thermo\endcsname{Thermodynamics, Bioenergetics, and Redox Calculations}
\expandafter\def\csname pname@f-membrane\endcsname{Membrane Potentials: Nernst, Goldman, and Cable Theory}
\expandafter\def\csname pname@f-transport\endcsname{Diffusion, Osmosis, Flow, and Water Potential}
\expandafter\def\csname pname@f-cardioresp\endcsname{Respiratory and Cardiovascular Calculations}
\expandafter\def\csname pname@f-renal\endcsname{Renal, Osmoregulatory, and Acid-Base Calculations}
\expandafter\def\csname pname@f-metabolic\endcsname{Metabolic Rate, Allometry, and Thermal Physiology}
\expandafter\def\csname pname@f-plantcalc\endcsname{Photosynthesis, Plant Water, and Growth Calculations}
\expandafter\def\csname pname@f-popgrowth\endcsname{Population Growth, Life Tables, and Demography}
\expandafter\def\csname pname@f-ecosystem\endcsname{Community, Diversity, and Ecosystem Formulas}
\expandafter\def\csname pname@f-epidemiology\endcsname{Epidemiology and Disease Dynamics}
\expandafter\def\csname pname@f-labcalc\endcsname{Laboratory, Molecular, and Bioinformatic Calculations}
\expandafter\def\csname pname@f-calculus\endcsname{Calculus, Differential Equations, and Modelling Shortcuts}
\expandafter\def\csname pname@f-allformulas\endcsname{Every Formula, Constant, and Normal Value}
\expandafter\def\csname pname@c-laws\endcsname{Named Laws, Rules, and Principles of Biology}
\expandafter\def\csname pname@c-pathways\endcsname{Metabolic Pathways and Cycles}
\expandafter\def\csname pname@c-molecules\endcsname{Molecules, Enzymes, Cofactors, and Vitamins}
\expandafter\def\csname pname@c-proteins\endcsname{Proteins, Receptors, Channels, and Transporters}
\expandafter\def\csname pname@c-genes\endcsname{Genes, Mutations, and Genetic Disorders}
\expandafter\def\csname pname@c-hormones\endcsname{Hormones and Signaling Molecules}
\expandafter\def\csname pname@c-experiments\endcsname{Model Organisms and Landmark Experiments}
\expandafter\def\csname pname@c-taxonomy\endcsname{Taxonomy: Every Phylum and Major Clade}
\expandafter\def\csname pname@c-anatomy\endcsname{Anatomy, Histology, and Normal Values}
\expandafter\def\csname pname@c-terminology\endcsname{Terminology, Roots, and the Timeline of Biology}
\newcommand{\partnameof}[1]{%
  \ifcsname pname@#1\endcsname\csname pname@#1\endcsname\else this book\fi}
\renewcommand{\pref}[1]{\@ifundefined{r@part:#1}{\emph{\partnameof{#1}}}{Part~\ref{part:#1}}}
\renewcommand{\prefn}[1]{\@ifundefined{r@part:#1}{\emph{\partnameof{#1}}}{\ref{part:#1}}}
\renewcommand{\nrrow}[2]{\@ifundefined{r@nr:#1}{#2}{#2\dotfill\pageref{nr:#1}}\\}
\makeatother

\begin{document}
\begin{titlepage}\thispagestyle{empty}
\vspace*{0.14\textheight}
\noindent{\color{tiercore}\rule{\textwidth}{5pt}}\\[3pt]
\noindent{\color{bioline}\rule{\textwidth}{1pt}}\\[26pt]
\noindent{\sffamily\bfseries\fontsize{11}{14}\selectfont\color{bionum}START HERE}\\[12pt]
\noindent\begin{minipage}{\textwidth}\raggedright
{\sffamily\bfseries\fontsize{30}{34}\selectfont\color{bioink}%
 \hyphenchar\font=-1 How to Learn All of This\par}\end{minipage}\\[16pt]
\noindent{\color{biogreen}\rule{\textwidth}{1.4pt}}\\[10pt]
\noindent\begin{minipage}{\textwidth}\raggedright
{\fontsize{10.5}{14}\selectfont\color{biogreen}%
Which file to open, the order to learn the topics in, what to memorise against what to
work out, and how to turn knowing a fact into recognising it in eight seconds under
time pressure.\par}\end{minipage}
\vfill
\noindent{\color{bioline}\rule{\textwidth}{1pt}}\\[3pt]
\noindent{\color{tiercore}\rule{\textwidth}{5pt}}
\end{titlepage}
\tableofcontents
\clearpage

\section{What is in these two zips and which file to open}

There are two zips. This one, \texttt{USABO-Biology-Concepts}, is the biology. The other,
\texttt{USABO-Biology-Formulas}, is the mathematics the same exams ask for. You need both,
but you do not need them at the same time, and the split exists so that you can revise
one without the other getting in the way.

\begin{keybox}[title={The four kinds of file}]
\begin{itemize}
\item \tm{The index of everything}, \texttt{00-Everything-Indexed.pdf}. Every fact,
      term, structure, pathway, organism, value and named result in the whole
      compendium, stated once with nothing explained, ordered exactly as the volumes
      are, with an alphabetical index at the back. Revise from this. Look things up in
      this.
\item \tm{Topic volumes}, in \texttt{01-Volumes/01-Core-Syllabus} and
      \texttt{01-Volumes/02-Beyond-the-Syllabus}. One topic each. Part one teaches it
      and works the methods. Part two is the bare memorization sheet for the same
      material. Learn from these.
\item \tm{Catalogue volumes}, in \texttt{01-Volumes/03-Reference-Catalogue}. Every named
      law, pathway, molecule, gene, hormone, experiment, taxon and word root, stated
      precisely. These are for looking up and for training recognition, not for reading
      end to end.
\item \tm{The formula volumes}, in the other zip. Derivations and worked calculations,
      area by area, plus their own index of every formula.
\end{itemize}
\end{keybox}

\begin{tipbox}[title={If you only do one thing}]
Open the topic volume for whatever you are weakest at. Read part one once, slowly, with
a pen. Then work from part two until you can reproduce every line of it from a blank
page. That single loop, repeated across topics in the order below, is the whole method.
Everything else in this file is refinement.
\end{tipbox}

\section{What the exams actually are}

\begin{keybox}[title={The competition, stage by stage}]
\begin{itemize}
\item \tm{The Open Exam.} Multiple choice, fifty minutes, taken at your school in
      February. Around a fifth of entrants advance. It rewards breadth and recall
      speed above all else.
\item \tm{The Semifinal Exam.} Two parts, around two hours, free response and data
      analysis alongside multiple choice. It rewards depth, mechanism, and the ability
      to reason from an unfamiliar figure. Twenty students advance to the National
      Finals.
\item \tm{The National Finals.} A residential week of theory and practical
      examinations. Four students are selected for the international team.
\item \tm{The IBO.} Two theory papers and four practical examinations. The practicals
      carry half the marks, which is why a practical volume exists in this compendium.
\end{itemize}
\end{keybox}

\begin{obsbox}[title={WORTH NOTICING: the official topic weightings}]
The IBO weightings, which the USABO follows closely, are: cell biology $20\%$, plant
anatomy and physiology $15\%$, animal anatomy and physiology $25\%$, ethology $5\%$,
genetics and evolution $20\%$, ecology $10\%$, biosystematics $5\%$. Two facts follow.
Animal physiology plus cell biology is $45\%$ of the exam, so those are where a week of
work pays the most. And ethology plus biosystematics is only $10\%$, but both are small
subjects that can be learned almost completely, so they are the cheapest marks in the
competition.
\end{obsbox}

\section{The order to learn in}

Topics are not independent. Learning them in dependency order means you never meet a
mechanism that rests on something you have not built yet. The order below is that
dependency order. It is not the order the volumes are numbered in, although the two are
close, because the numbering also has to keep related volumes adjacent on disk.

\begin{keybox}[title={The dependency order for the core syllabus}]
\begin{enumerate}
\item \tm{Foundations, Scale, and the Language of Biology.} Units, magnitudes,
      experimental design, graph reading, microscopy arithmetic, and word roots. Short,
      and every later volume assumes it. Do not skip it because it looks easy: the
      unit conversions and the data figure conventions here are tested directly.
\item \tm{The Chemistry of Life}, then \tm{Biological Macromolecules}, then
      \tm{Enzymes, Energetics, and Metabolic Control}. This block is the substrate of
      everything. Free energy, pH, bonding and protein structure are used in every
      volume after it, and enzyme kinetics is used in the formula zip as well.
\item \tm{Cell Structure and the Organelles}, then \tm{Membranes and Membrane
      Transport}. Transport before metabolism, because chemiosmosis is a transport
      argument.
\item \tm{Cellular Respiration and Fermentation}, then \tm{Photosynthesis}. Learn them
      in that order and in the same week: photosynthesis is respiration's machinery run
      backwards, and learning them together halves the work.
\item \tm{Cell Signaling}, then \tm{The Cytoskeleton, Motility, and the Extracellular
      Matrix}, then \tm{The Cell Cycle, Mitosis, and Meiosis}, then \tm{Cancer,
      Apoptosis, and the Control of Cell Number}. Meiosis is the hinge into genetics,
      so do not leave it half learned.
\item \tm{DNA Replication}, then \tm{Transcription and the Genetic Code}, then
      \tm{Translation and Protein Targeting}, then \tm{Regulation of Gene Expression and
      Epigenetics}, then \tm{Biotechnology and Molecular Laboratory Methods}. Strictly
      in this order. Biotechnology last, because every technique in it is an
      application of a mechanism from the four volumes before it.
\item \tm{Mendelian Genetics}, then \tm{Linkage and Mapping}, then \tm{Population
      Genetics and Hardy-Weinberg}. Take the formula zip's genetics volumes alongside
      these three, not afterwards: the mathematics and the biology are the same
      material seen from two sides.
\item \tm{Evolutionary Mechanisms}, then \tm{Speciation and Macroevolution}, then
      \tm{Phylogenetics and Biosystematics}. Population genetics is the prerequisite
      for all three.
\item \tm{Bacteria and Archaea}, \tm{Viruses}, \tm{Protists}, \tm{Fungi}. Diversity is
      pure recall, so it is the block to learn when you are tired, and the block to
      revisit most often. Phylogenetics first, so that the groups hang on a tree
      instead of on a list.
\item \tm{Plant Diversity and Reproduction}, \tm{Plant Anatomy}, \tm{Plant Transport
      and Nutrition}, \tm{Plant Hormones and Development}. Anatomy before transport,
      because every transport argument is about a named tissue.
\item \tm{Animal Diversity and Body Plans}, then the physiology volumes. Within
      physiology the order that minimises repeated work is: \tm{Tissues, Skeleton, and
      Muscle}, \tm{Nutrition and Digestion}, \tm{Gas Exchange}, \tm{Circulation},
      \tm{Immunology}, \tm{Osmoregulation and Thermoregulation}, \tm{Nervous System},
      \tm{Sensory Biology}, \tm{Endocrine System}, \tm{Reproduction and Development}.
      Gas exchange immediately before circulation, and the nervous system immediately
      before the endocrine system, because each pair shares its control logic.
\item \tm{Ethology and Animal Behavior}, then \tm{Population and Community Ecology},
      then \tm{Ecosystems and Conservation}. Behaviour needs evolution; ecology needs
      behaviour and population genetics.
\item \tm{Exam Craft} last, and then again in the final week.
\end{enumerate}
\end{keybox}

\begin{warnbox}[title={TRAP: reading in folder order}]
The folder order is not the learning order. The largest single mistake students make
with a compendium like this is to start at volume one and read forwards, which puts
diversity and plant anatomy before the molecular biology they are described with, and
puts exam craft at the end when it should shape how you read everything else. Use the
list above, and read the exam craft volume in week one as well as in the final week.
\end{warnbox}

\section{What to memorise and what to work out}

Biology rewards memory more than any other olympiad subject, and that fact is used as
an excuse for learning nothing but lists. The distinction below is worth more than any
amount of extra reading.

\begin{keybox}[title={Memorise these, cold, with no thinking time}]
\begin{itemize}
\item Every pathway's inputs, outputs, location, and net yield. Glycolysis, the citric
      acid cycle, the electron transport chain, the Calvin cycle, beta oxidation, the
      urea cycle.
\item The twenty amino acids by class, the genetic code's start and stop codons, the
      base pairing rules, and the four macromolecule linkages.
\item Every organelle and its function, every plant and animal tissue, every human
      organ system's components.
\item The hormones: source, chemical class, target, action. There are about fifty that
      matter and they are pure recall.
\item Normal physiological values. Blood pH $7.35$ to $7.45$, plasma $\mathrm{Na^+}$
      $135$ to $145\ \mathrm{mM}$, resting potential $-70\ \mathrm{mV}$, cardiac output
      $5\ \mathrm{L\,min^{-1}}$. Questions are built on these.
\item The taxonomic groups and their defining synapomorphies.
\item The named experiments: who, what design, what result, what it proved.
\end{itemize}
\end{keybox}

\begin{keybox}[title={Work these out, every time, from a principle}]
\begin{itemize}
\item Anything with a direction: water movement, ion movement, a shift in an
      equilibrium, the sign of a free energy change. Memorising directions produces
      confident wrong answers; deriving them from a gradient takes four seconds.
\item Genetic ratios. Learn the mechanism of epistasis and derive the $9{:}3{:}4$, do
      not memorise a table of modified ratios.
\item Any feedback loop's response to a perturbation. Find the sensor, the effector and
      the sign, and the answer follows.
\item Anything the formula zip covers. A formula you can reconstruct is worth three you
      have memorised, because the exam perturbs the conditions.
\end{itemize}
\end{keybox}

\begin{tipbox}[title={TECHNIQUE: the recognition drill}]
The slow step in a biology exam is not recall, it is recognition: seeing which of four
hundred mechanisms a question is about. Train it directly. Open the index of everything
at a random page, read one entry name only, and say aloud what it is, where it sits, and
what question would be built on it. Five minutes of this is worth an hour of rereading,
because it trains the retrieval direction the exam actually uses.
\end{tipbox}

\section{How to use the memorization sheets}

\begin{keybox}[title={The loop}]
\begin{enumerate}
\item Read part one of a volume once, with a pen, doing the worked calculations
      yourself before reading the solution.
\item Close it. Write down everything you remember, on blank paper, in any order.
\item Open part two, the memorization sheet, and mark what you missed.
\item Work only the missed entries the next day, then the whole sheet again three days
      later, then a week later.
\end{enumerate}
That is spaced repetition with the interval chosen for you. The only part that matters
is step two: production from a blank page, not recognition of a page you are looking at.
\end{keybox}

\begin{warnbox}[title={TRAP: highlighting feels like learning}]
Rereading and highlighting produce a strong feeling of fluency and almost no retention.
Every study that has compared them with active recall has found the same thing. If you
are not producing material from an empty page, you are not learning it, however
familiar it feels.
\end{warnbox}

\section{Twelve weeks}

The plan below assumes about ten hours a week. Each week ends with one hour of mixed
recall drawn from every week so far, which is the part people drop and the part that
does the work.

\begin{keybox}[title={The twelve week plan}]
\begin{enumerate}
\item Foundations, exam craft, and the chemistry of life. Set up the recall loop.
\item Macromolecules and enzymes. Formula zip: units, and enzyme kinetics.
\item Cell structure, membranes, transport. Formula zip: diffusion and membrane
      potentials.
\item Respiration and photosynthesis together. Formula zip: thermodynamics and redox.
\item Signalling, cytoskeleton, cell cycle, cancer.
\item Replication, transcription, translation, gene regulation.
\item Biotechnology, then Mendelian genetics and linkage. Formula zip: probability,
      Mendelian and linkage calculations.
\item Population genetics and evolution. Formula zip: Hardy-Weinberg and selection.
\item Speciation, phylogenetics, and the whole diversity block.
\item All four plant volumes. Formula zip: plant calculations.
\item Animal physiology, first half: tissues and muscle, digestion, gas exchange,
      circulation, immunology. Formula zip: cardiorespiratory calculations.
\item Animal physiology, second half: osmoregulation, nervous, sensory, endocrine,
      reproduction and development. Then ethology and ecology. Formula zip: renal,
      metabolic rate, population and community formulas.
\end{enumerate}
\end{keybox}

\begin{tipbox}[title={The four week compression}]
If four weeks is what you have, the ordering changes because breadth now beats depth.
\begin{enumerate}
\item Read only the memorization sheets, in the dependency order above, at the rate of
      about twelve a day, and mark every entry you do not already know.
\item Week two: read part one of the volumes covering the marked entries only.
\item Week three: cell biology and animal physiology in full, because they are $45\%$
      of the exam. Then the formula zip's index of every formula, straight through.
\item Week four: the index of everything, front to back, twice, plus exam craft and one
      timed paper. Nothing new after day three of this week.
\end{enumerate}
\end{tipbox}

\section{The final week}

\begin{keybox}[title={What to do and what to stop doing}]
\begin{itemize}
\item Stop reading part one of anything. New material this week displaces retrievable
      material.
\item Work the index of everything front to back. It exists for this week.
\item Work the index of every formula the same way, and rebuild each equation from its
      principle rather than reading it.
\item Re-read the exam craft volume's sections on multiple choice technique, negation
      traps, and time allocation.
\item Spend one session on the diversity and taxonomy sheets. They decay fastest and
      they are the cheapest marks to recover.
\item Sleep. Recall is the first thing sleep deprivation takes, and recall is the whole
      exam.
\end{itemize}
\end{keybox}
\end{document}
